\section{Discussion}

The central result is a representational constraint rather than a new claim
about hippocampal capacity. In this minimal architecture, a CA3-like retrieval
key selects associative content, plastic CA3--CA1 weights express that content
as a CA1 state, and a fixed decoder determines whether the state remains usable
at the output. A coordinate permutation preserves the instructed vector's
values yet disrupts reconstruction when the decoder is not transformed with
it. The matched-decoder rescue makes the logic exact: rapid plasticity must
preserve, or co-adapt with, the representational coordinates assumed by
recipient circuits. This motivates viewing CA1 not only as a locus of
experience-dependent activity, but also as an adaptive interface that places
retrieved content within a structured neural space. Such a space could allow
successive experiences to remain decodable even while individual fields and
participating neurons change.

This proposal complements rather than replaces existing accounts of
hippocampal memory. Autoassociative CA3 theories explain how a stored state can
be selected from a partial cue, whereas BTSP models explain how temporally broad
instructive events can rapidly alter CA1 synapses
\citep{nakazawa2002,bittner2017,wu2025}. The present model abstracts selection
as a non-recurrent CA3-like key and asks what must additionally hold when the
selected state is expressed at hippocampal output. The relationship to
neural-manifold work is conceptual: those studies concern motor cortical
learning, not hippocampal memory, but similarly show that an unchanged readout
assigns meaning only to particular population coordinates
\citep{sadtler2014,golub2018}. Readout compatibility is therefore a general
systems constraint instantiated here in a hippocampal circuit motif. The
simulations do not test retention across days or years; long-term decodability
is a hypothesis about what preserving this relation might enable, not an
observed outcome.

Pretraining and rapid plasticity have distinct computational roles in this
proposal. Pretraining establishes a sparse representational basis together with
the decoder that gives its coordinates meaning; rapid CA3--CA1 learning then
places a new, key-addressable experience within that basis. This resembles
feature reuse after pretraining in deep networks \citep{yosinski2014}, but the
analogy is functional rather than mechanistic. Likewise, the covariance matrix
adaptation evolution strategy (CMA-ES) used during preliminary model
development was an engineering procedure for locating a workable parameter
regime, not a proposed neural learning mechanism. In a biological circuit, the
basis could reflect development, prior learning, and ongoing
entorhinal--hippocampal organization rather than a separate optimization phase.
The frozen decoder is an assay of compatibility, not a claim that biological
recipient pathways are immutable. The broader prediction is that changes in a
CA1 code and its downstream readouts must remain coordinated.

The cue-track simulations connect readout compatibility to context-dependent
reconfiguration. Exchanging cue identities produced abrupt changes in recalled
CA1 tuning under both plasticity rules, whereas matched no-swap controls
remained stable at the same lap boundaries. Event alignment showed that most
of the change occurred on the first exchanged-cue lap and that consecutive-lap
similarity largely recovered on the next lap. Population-map similarity formed
alternating blocks: representations were similar within a cue assignment and
again when that assignment recurred, rather than drifting monotonically across
the experiment. After learning, both rules produced continuous mixtures of
spatially stable and cue-modulated units, with broad unit-level endpoint
similarities. This behavior is qualitatively consistent with cue-sensitive
field and rate changes and with evidence that EC-directed BTSP can form, shift,
and stabilize CA1 fields around behaviorally relevant features
\citep{leutgeb2005,grienberger2022,madar2025,vaidya2025}. On this
interpretation, instructed plasticity acts as an update operation on a spatial
cognitive map: it can insert, weaken, or reposition fields as relevant cues
change. In the present two-context task, the block structure is better
described as cue-dependent re-expression and updating than as the formation of
an unconstrained new map at every transition. The simulations are not a
quantitative model of biological remapping, however, and the unit distributions
were not compared with recorded field statistics. Their narrower implication
is that reconfiguration should be evaluated not only by how much CA1 activity
changes, but also by whether the new population state remains interpretable at
the output.

Comparing the two update rules clarifies why the form of rapid plasticity
matters. The instructive-driven rule produced stronger one-shot reconstruction and a larger
drop in tuning similarity at cue exchanges, followed by a more gradual recovery
over the next laps. The error-driven rule performed a bounded residual
correction and recovered most consecutive-lap stability by the following lap;
after repeated track experience it also produced greater spatial stability and
much higher clean position accuracy. These
differences should not be reduced to a single ranking. The instructive-driven
rule is closer in spirit to an instructed write, whereas the error-driven rule
adds a computational mechanism
for correcting under- and overexpressed CA1 components. Because both rules used
one shared parameter regime, their comparison reveals consequences of the
update signal under matched conditions rather than the best performance each
could achieve. A stronger model comparison would optimize each rule
independently and vary memory load, cue overlap, and learning rate.

The degradation experiments distinguish robustness to partial input from
mechanistic pattern completion. With the normal key map, recall remained
informative after substantial input loss, and selective removal of MEC-like or
LEC-like activity preferentially affected the corresponding readout. Yet the
CA3-like population has no recurrent connections or attractor settling.
Robustness arises from a distributed feedforward input-to-key map followed by a
learned key-to-content association. Degraded-cue recall alone therefore does
not identify a recurrent completion mechanism. A biological claim about CA3
pattern completion would require explicit recurrent dynamics and comparison
with interventions that alter those dynamics \citep{nakazawa2002,li2024}.

The dense common-key control also has a deliberately narrow interpretation. It
forces all CA3-like units to receive the same EC average and thus destroys
item-specific key structure while changing connection density and sparsity. Its
failure shows that a non-specific common key cannot support these associations;
it does not show that the selected sparse wiring is uniquely necessary or
optimal. Likewise, the MEC/LEC dissociation is partly expected from the
constructed input and readout partitions. Its value is as an internal
validation that the shared CA1 basis can retain access to an intact component,
not as evidence for strict anatomical modularity. Sparsity-matched controls
would be needed to separate the effects of connection density, dimensionality,
and key discriminability.

Several additional limitations define the scope of the work. The model omits
dentate gyrus, recurrent CA3 dynamics, attractor competition, inhibition,
consolidation, and explicit long-term drift of neurons or readouts. MEC and LEC
are simplified input halves despite mixed selectivity in both regions. The CA1
basis is obtained from a shallow autoencoder rather than a biologically
specified developmental process, and the plasticity equations are bounded
abstractions rather than temporal or molecular models of BTSP. The
compatibility manipulation is deliberately synthetic, and its matched-decoder
rescue is an equivariance control rather than independent biological evidence.
Finally, the track contains only two cue identities at one memory load;
capacity, interference, and scaling across multiple cue-pair sequences remain
unresolved.

These limitations suggest direct next tests. Model performance should depend
more strongly on compatibility between an instructed CA1 state and its output
basis than on the magnitude of CA1 change alone. Degraded-cue recall should
track retrieval-key discriminability across sparsity-matched controls, and
increasing the number of cue-pair sequences should eventually reveal an
interference-dependent capacity limit. Experimentally, the broad prediction is
that learning-related CA1 reconfiguration can remain behaviorally useful when
recipient populations preserve or acquire an appropriate readout, whereas an
equally large but readout-incompatible change should impair retrieval.

Within these bounds, the model offers a compact proposal linking a previously
established representational basis, rapid instructed plasticity,
retrieval-key-based selection, and downstream readability. CA1 activity can
change as a cognitive map is updated yet remain useful if retrieved content is
expressed in coordinates that recipient circuits can interpret. Readout
compatibility is therefore a useful axis for evaluating models of rapid
hippocampal learning and for asking how a changing hippocampal code might
support decodable experience over longer timescales.
